\documentclass[11pt]{article}

\newcommand{\cnum}{Math 245B}
\newcommand{\ced}{Winter 2026}
\newcommand{\ctitle}[4]{\title{\vspace{-0.5in}\cnum, \ced\\Assignment #1: #2}\author{\vspace{-0.35in}\\#3, #4}}
\usepackage{enumitem}
\usepackage[usenames,dvipsnames,svgnames,table,hyperref]{xcolor}
\usepackage{amsmath, amsfonts}
\usepackage{graphicx} % Required for inserting images
\usepackage{float}
\usepackage{listings}
\usepackage{xcolor}
\usepackage{hyperref}
\usepackage{comment}

\renewcommand*{\theenumi}{\alph{enumi}}
\renewcommand*\labelenumi{(\theenumi)}
\renewcommand*{\theenumii}{\roman{enumii}}
\renewcommand*\labelenumii{\theenumii.}

\definecolor{codegreen}{rgb}{0,0.6,0}
\definecolor{codegray}{rgb}{0.5,0.5,0.5}
\definecolor{codepurple}{rgb}{0.58,0,0.82}
\definecolor{backcolour}{rgb}{0.95,0.95,0.92}
\definecolor{header}{HTML}{205459}
\definecolor{solution}{HTML}{204d52}

\newcommand{\solution}[1]{{{\textcolor{header}{Solution:} \textcolor{solution}{#1}}}}
\setlist[enumerate]{label=(\alph*)}

\lstdefinestyle{mystyle}{
    backgroundcolor=\color{backcolour},
    commentstyle=\color{codegreen},
    keywordstyle=\color{magenta},
    numberstyle=\tiny\color{codegray},
    stringstyle=\color{codepurple},
    basicstyle=\ttfamily\footnotesize,
    breakatwhitespace=false,
    breaklines=true,
    captionpos=b,
    keepspaces=true,
    numbers=left,
    numbersep=5pt,
    showspaces=false,
    showstringspaces=false,
    showtabs=false,
    tabsize=2
}


\begin{document}
\ctitle{\#1}{}{Asher Christian}{006-150-286}
\date{}
\maketitle
\section{intro}
Stochastic = random, Proces = sequence of things
SP = sequence of random things
Independent identically distributed random variables, $X_1,X_2,\dots,x_n,\dots$
Ex Partial sum of IID Rvs:
\[
S_0 = 0, S_1 = X_1, S_2 = X_1 + X_2, \dots, S_n = \sum_{i=1}^{n}X_ii, \dots
\] 
\[
    S_n = S_{n-1} + x_n
\] 
if $X_n \in \{0,1\}$ then  $S_n = S_{n-1}$ or  $S_n = S_{n-1} + 1$ we can say
 \[
     \frac{S_n}{n} \rightarrow \mathbb{E}[X_1]
\] 
by SLLN almost sure
\[
    \frac{S_n - n \mathbb{E}[x_1]}{\sqrt{n}\sqrt{Var(X_1)}} \rightarrow^{d} N(0,1)
\] 
"simple model of dependence Given data at present, the future is independent of the past"\\
Andrei Markov: Markov chain in 1906
\[
X_0,x_1,X_2,\dots,X_n,\dots
\] 
\[
    (*)  \;\;\;  \mathbb{P}(X_{n+1} = j|X_n = i_n, X_{n-1} = i_{n-1}, \dots, X_1 = i_1, x_0 = i_0) = \mathbb{P}(X_{n+1} = j| X_n = i_n)
\] 
Any SP $\{X_n : n \ge 0\}$ satisfying  $(*)$ for all  $n \ge 0$ all  $j,i,i_{n-1},\dots,i_0$is called a Markov chain.
"Time-homogenous MC": RHS of $(*)$ does not depend on time so it is equal to  $p(i,j)$\\
$p(i,j)$ is called the transition probability of $X_n$
 \[
\mathbb{P}(X_2 = j | X_1 = i ) = p(i,j)
\] 
\[
\mathbb{P}(X_2 = j | X_0 = i) = \sum_{k}^{} \mathbb{P}(X_2 = j, X_1 = k | X_0 = i)
\] 
\[
 = \sum_{k}^{} \mathbb{P}(X_2 = j | X_1 = k, X_0 = i) \mathbb{P}(X_1 = k | X_0 =i) = \sum_{k}^{}p(k,j)p(i,k)
\] 
Conditions:
\[
\begin{cases}
    0 \le p(i,j) \le 1 \\
    \sum_{j}^{}p(i,j) = 1
\end{cases}
\]
Example (Gambling):\\
Fair game equal probablity to win $\$1$ or lose  $\$1$. You continuou playing until
\begin{enumerate}
    \item you lose all your money
    \item you gained $N$ dollars
\end{enumerate}
Predict money $X_n$ after the  $n$ round.\\
$p(i,i+1) = \frac{1}{2} = p(i,i-1)$ given $1 \le i \le N-1$\\
$p(0,0) = 1 = p(N,N)$\\
Finite state MC. $\{0,1,2,\dots,N\}$ only  $(n+1)^2$ transition probabilities.\\
What if $X_0 \sim \mu_0$ $ \mathbb{P}(X_0 = i) = \mu_0(i)$
$\mu_1(j) = \mathbb{P}(X_1 = j | X_0) = \sum_{i}^{}\mu_0(i) \mathbb{P}(X_1 = j| X_0 = i) =\sum_{i}^{}\mu_0(i) p(i,j)$
\[
\mu = (\mu_0(1),\mu_0(2),\dots,\mu_0(N))
\] 
\[
    \mu_1(j) = \mu_0 * P(j-\text{th column}) = \mu_0 \circ P=  \mu_0 P
\] 
$P_{i,j} = p(i,j)$

\section{Gene mutation}
DNA sequence $\{A,T,G,C\}^{d}$ 
notation  $d_H(D_1,D_2) = \sum_{i=1}^{d}1_{D_1(i) \ne D_2(i)}$.
$X_n = d_H(X_n, D_0)$ $X_n$ lies in  $\{0,1,2,\dots,d\}$
$p(d,d-1) = 1, p(0,1) = 1$
\[
p(i,i+1) = ? \qquad p(i,i-1) = ? \qquad 1 \le i \le d-1
\] 
\[
    p(i,i-1) = \mathbb{P}(d_H(D_{N+1},D_0) = i-1 | D_H(D_n, D_0) = i) = \frac{i}{d}
\] 
since there are $i$ elements that are wrong and so their proportion is $\frac{i}{d}$
\[
    p(i,i+1) = \frac{d-i}{d}
\] 

\section{Branching Process}
$p_k : k = 0,1,\dots$
\[
\sum_{k=0}^{\infty}p_k = 1 \qquad p_k \ge 0
\] 
where $p_k = \mathbb{P}(\text{each parent gives birth to $k$ offsprint})$ each parent gives birth independently according to this model.
$X_n = \text{total number of children at generation $n$ }$ 
$z \sim (p_k)_{k=1}^{\infty}$
\[
    X_{n+1} = \sum_{k=1}^{X_n}Z_{n,k}
\] 
we have $Z_{n,k}$ is the number of kids the  $k$-th kid has in the  $n$-th generation
 \[
     p(i,j) = \mathbb{P}(X_{n+1} = j | X_n = i) = \mathbb{P}(\sum_{k=1}^{i}Z_{n,k} = j)
\] 

\end{document}
